\chapter{GIỚI THIỆU TỔNG QUAN}
\label{ch:gioi_thieu}

\section{Đặt vấn đề}

Xe tự hành và các hệ thống hỗ trợ lái tiên tiến (ADAS) phải hiểu được môi trường giao thông xung
quanh trước khi có thể ra bất kỳ quyết định điều khiển nào. Trong số các thành phần của bài toán
nhận thức đó, phát hiện biển báo giao thông giữ một vị trí đặc biệt: biển báo là kênh giao tiếp
chính thức giữa hạ tầng đường bộ và người lái, và nội dung của nó ràng buộc trực tiếp hành vi hợp
pháp của phương tiện. Một hệ thống bỏ sót biển giới hạn tốc độ hoặc đọc nhầm đèn tín hiệu không chỉ
gây khó chịu mà có thể dẫn tới hậu quả về an toàn và pháp lý.

Về mặt kỹ thuật, đây là bài toán phát hiện đối tượng: với mỗi ảnh đầu vào, mô hình phải đồng thời
\emph{định vị} từng biển báo bằng một hộp bao và \emph{phân loại} nó vào đúng chủng loại. Các bộ
phát hiện một giai đoạn thuộc họ YOLO~\cite{redmon2016yolo,jocher2023ultralytics} hiện là lựa chọn
mặc định cho ứng dụng thời gian thực nhờ chỉ cần một lượt truyền xuôi.

Tuy nhiên, đạt được một chỉ số mAP cao trên tập kiểm tra mới chỉ là một nửa bài toán. Nửa còn lại
--- và cũng là nửa quyết định việc hệ thống có triển khai được hay không --- là ràng buộc tài
nguyên. Máy tính gắn trên xe không có card đồ hoạ máy chủ; nó có ngân sách bộ nhớ, băng thông và
điện năng rất hạn hẹp. Do đó nhu cầu \emph{nén mô hình} trở thành trung tâm: làm sao để mô hình nhỏ
hơn và nhanh hơn mà không đánh mất năng lực phát hiện.

Hai họ kỹ thuật nén phổ biến nhất đi theo hai hướng khác nhau. \textbf{Chưng cất tri
thức}~\cite{hinton2015distillation} tác động ở khâu huấn luyện: một mô hình học trò nhỏ được dạy
bắt chước phân phối đầu ra mềm của một mô hình thầy lớn hơn, với kỳ vọng thu được năng lực vượt quá
mức mà kích thước của nó cho phép. \textbf{Lượng tử hoá sau huấn
luyện}~\cite{jacob2018quantization} tác động ở khâu triển khai: thay biểu diễn dấu phẩy động bằng
biểu diễn số nguyên độ rộng bit thấp hơn, không cần huấn luyện lại.

Cả hai kỹ thuật đều được trình bày rộng rãi kèm những con số hấp dẫn. Nhưng những con số đó gần như
luôn được báo cáo dưới dạng \emph{mAP tổng} --- và đó chính là điểm mà đề tài này đặt nghi vấn.

Các nút thắt chính mà đề tài phải đối mặt:

\begin{enumerate}[leftmargin=1.5cm]
\item \textbf{Ràng buộc kép giữa độ chính xác và chi phí triển khai.} Mô hình phải đủ nhẹ để chạy
trên phần cứng nhúng nhưng vẫn phải đủ chính xác để không bỏ sót biển báo.

\item \textbf{mAP tổng có thể là một chỉ số gây hiểu nhầm.} Nó lấy trung bình trên toàn bộ vật thể
bất kể kích thước. Nếu phần lớn vật thể trong tập kiểm tra là những trường hợp dễ, một suy giảm
nghiêm trọng ở nhóm khó có thể bị pha loãng gần như hoàn toàn.

\item \textbf{Vật thể nhỏ.} Biển báo ở xa chỉ chiếm vài chục điểm ảnh trong khung hình, và đây lại
chính là trường hợp quyết định thời gian phản ứng của xe. Phát hiện vật thể nhỏ vốn là điểm yếu cố
hữu của các bộ phát hiện hiện đại~\cite{tong2020smallobject}.

\item \textbf{Dữ liệu quy mô vừa và mất cân bằng.} Bộ dữ liệu sử dụng chỉ có 6.012 hộp bao cho 15
lớp, với tỉ số mất cân bằng giữa lớp phổ biến nhất và lớp hiếm nhất lên tới 35,77
lần~\cite{johnson2019imbalance}.

\item \textbf{Nguy cơ rò rỉ dữ liệu giữa các tập.} Bộ dữ liệu công khai thường được ghép từ nhiều
nguồn và có thể chứa ảnh trùng lặp giữa tập huấn luyện và tập kiểm tra, khiến chỉ số bị thổi
phồng~\cite{barz2020duplicates}.
\end{enumerate}

\noindent\textbf{Câu hỏi nghiên cứu 1.1} \textit{Chưng cất tri thức từ một mô hình thầy lớn hơn có
thực sự giúp mô hình học trò nhỏ tốt hơn so với chính nó khi được huấn luyện thông thường, trong
điều kiện bộ dữ liệu quy mô vừa?}

\noindent\textbf{Câu hỏi nghiên cứu 1.2} \textit{Lượng tử hoá sau huấn luyện đánh đổi những gì
giữa dung lượng, tốc độ và độ chính xác, và cái giá đó phân bố ra sao trên các nhóm vật thể có kích
thước khác nhau?}

\noindent\textbf{Câu hỏi nghiên cứu 1.3} \textit{Chỉ số mAP tổng phản ánh đến mức nào năng lực thực
sự của mô hình, và đặc điểm nào của bộ dữ liệu khiến chỉ số đó có nguy cơ gây hiểu nhầm?}

\section{Mục tiêu đề tài}

Đề tài hướng đến các mục tiêu chính sau:

\begin{itemize}[leftmargin=1.5cm]
\item Xây dựng một quy trình thực nghiệm tái lập được, trong đó cả năm cấu hình khảo sát được huấn
luyện và đánh giá trong \emph{cùng một lượt chạy}, trên cùng một tập kiểm tra và cùng một bộ tính
độ đo, để loại bỏ mọi khác biệt do môi trường.

\item Đánh giá bằng pycocotools trên định dạng COCO nhằm thu được $AP$ tách riêng theo nhóm kích
thước vật thể, thay vì chỉ báo cáo mAP tổng.

\item Khai báo trước các ngưỡng chấp nhận cho từng giả thuyết, và báo cáo đầy đủ cả những ngưỡng
không đạt.

\item Phân tích khám phá dữ liệu ở mức đủ sâu để giải thích được kết quả, bao gồm kiểm toán rò rỉ
giữa các tập.

\item Báo cáo trung thực các kết quả đi ngược giả thuyết ban đầu, kèm phân tích nguyên nhân và các
yếu tố gây nhiễu chưa kiểm soát được.
\end{itemize}

\section{Đối tượng và phạm vi nghiên cứu}

\textbf{Đối tượng nghiên cứu:} hiệu quả của hai kỹ thuật nén mô hình --- chưng cất tri thức và
lượng tử hoá sau huấn luyện --- áp dụng cho bộ phát hiện biển báo giao thông, xét đồng thời trên
ba trục độ chính xác, tốc độ và dung lượng, trong đó độ chính xác được tách theo nhóm kích thước
vật thể.

\textbf{Phạm vi nghiên cứu:} đề tài giới hạn ở bộ dữ liệu \textit{pkdarabi/cardetection} với 15 lớp
gốc, giữ nguyên không ánh xạ lại nhãn. Mô hình thầy là YOLO26s, mô hình học trò là YOLO26n, cả hai
huấn luyện 50 chu kỳ ở kích thước ảnh 640. Lượng tử hoá được thực hiện qua đường xuất ONNX.

Đề tài \emph{không} bao gồm: thu thập và gán nhãn dữ liệu mới, tìm kiếm siêu tham số tự động, quét
hệ số chưng cất, khảo sát các kiến trúc phát hiện khác, và triển khai thực tế trên phần cứng nhúng.
Những giới hạn này được nêu rõ ngay từ đầu vì chúng ảnh hưởng trực tiếp tới cách diễn giải kết quả
ở Chương~\ref{ch:ket_qua}.

\section{Đóng góp của đề tài}

\begin{enumerate}[leftmargin=1.5cm]
\item \textbf{Một thiết kế thực nghiệm có đối chứng đúng nghĩa.} Mô hình học trò được huấn luyện
hai lần với cùng số chu kỳ, cùng hạt ngẫu nhiên, cùng kích thước ảnh --- một lần thông thường và
một lần có chưng cất --- nên chênh lệch quan sát được quy về đúng biến số cần khảo sát. Kèm theo đó
là năm ngưỡng chấp nhận khai báo trước khi chạy.

\item \textbf{Bằng chứng định lượng cho thấy mAP tổng che giấu chi phí thật của việc nén mô hình.}
Lượng tử hoá INT8 làm mAP@0,5 giảm 0,0038 nhưng làm $AP$ trên vật thể nhỏ giảm 0,0247 --- gấp 6,5
lần. Đây là đóng góp chính của đề tài và nó có hệ quả thực tiễn trực tiếp: mọi quyết định chọn cấu
hình triển khai chỉ dựa trên mAP tổng đều có nguy cơ sai.

\item \textbf{Một kết quả âm được báo cáo đầy đủ.} Chưng cất tri thức làm giảm độ chính xác so với
đối chứng. Kết quả này được giữ lại kèm bằng chứng cho thấy cơ chế chưng cất đã thực sự hoạt động,
và kèm phân tích các điều kiện khiến kỹ thuật không phát huy tác dụng.

\item \textbf{Một phân tích dữ liệu chỉ ra hai vấn đề của bộ dữ liệu công khai được dùng rộng rãi:}
38,68\% số ảnh là ảnh cắt cận cảnh kiểu GTSRB~\cite{stallkamp2012gtsrb} khiến phân bố kích thước bị
lưỡng cực, và 10,2\% tập kiểm tra có bản sao trong tập huấn luyện.
\end{enumerate}

\section{Bố cục báo cáo}

Chương~\ref{ch:gioi_thieu} giới thiệu tổng quan về bài toán, mục tiêu và phạm vi.
Chương~\ref{ch:co_so} trình bày cơ sở lý thuyết về phát hiện đối tượng một giai đoạn, hệ thống độ
đo và hai kỹ thuật nén mô hình. Chương~\ref{ch:phuong_phap} mô tả bộ dữ liệu, kết quả phân tích
khám phá và thiết kế thực nghiệm. Chương~\ref{ch:ket_qua} trình bày và bàn luận kết quả, bao gồm
phần phân tích lỗi và các kết quả trái với kỳ vọng. Chương~\ref{ch:tong_ket} tổng kết, nêu hạn chế
và hướng phát triển.
